\documentclass[11pt]{amsart}

\usepackage{amsmath, amssymb, amsthm}
\usepackage{mathtools}
\usepackage{thmtools}
\usepackage{enumitem}
\usepackage{tikz-cd}
\usepackage{hyperref}
\usepackage[nameinlink]{cleveref}

% % Theorem environments
\declaretheorem[numberwithin=section]{theorem}
\declaretheorem[style=plain, sibling=theorem]{lemma, proposition, corollary}
\declaretheorem[style=definition]{definition, example, condition}
\declaretheorem[style=remark]{remark, note, observation, todo}


\title{April 17, 2026}


\begin{document}
\maketitle

I had a regular team seminar in this morning.
I expected that I have a talk today, but (fortunately) Dr. Seonhwa Kim gave a talk about his recent reserch.
It was about ED-rigidicy for nonconvex polyhedra.
I record short review for it for the future.

Suppose that we have a 3-dimensional polyhedron in Eulidiean space and we know lengths of all the edges and all the dihedral angles.
In this situation, the question is that, whether this polyhedron is geometrically rigid or not, i.e., every polyhedron with the same edge lengths and dihedral angles are congruent.
It is well-known that the answer is yes for convex polyhedra, but it is not known for nonconvex ones.
Seonhwa Kim proved that the answer is also yes for nonconvex ones adapting the notion of rigid vertices.

The vertex $v$ is called \emph{rigid} if $\mathrm{nt}(v) \leq 3$ where $\mathrm{nt}(v)$ is the number of triangles among neighboring faces of $v$.
The number 3 is related to the fact that if we know all the lengths of the link of a vertex $v$, then we know all the facial angles if $\mathrm{nt}(v) \leq 3$.
There is a proposition which states that There exists a rigid vertex for any polyhedra.

He proved the theorem by using induction on the nubmer of vertices of a polyhedron.
At a rigid vertex $v$, by removing the ``cone'' of $v$ from the given polyhedron we can remove one vertex from it.
The point is that, we know all the edge lengths and dihedral angles of the link of $v$, and that the remaining polyhedron is also in ED condition.
The caution is that, we may add edges and faces when we do remove, and how they are added is depending on the geometric conditon of the polyhedron, not the combinatorial one.
In other words, the statement that we have a polyhedron after removing the cone is not free, it is a corollary of the\emph{Steinitz theorem}.

\par\vspace{1em}

After the talk, I went to the hairshop for cutting my hair. 
Then I first reviewed updated papers from arXiv, and then I checked the generated feedback file that codex generated.
It reports that there were lots of typos, wrong article usage (still), and non idiomatic phrases.
I again think that I should a find a way to become accustomed to idioms, I will consult a plan with GPT in this weekend.

\par\vspace{1em}

After then I tried to my own job .. from 6 pm.





\bibliographystyle{amsalpha}
% \nocite{*}
\bibliography{bibliography}
\end{document}
